\worldchapter{Vorwort}{Foreword}
\label{ch:vorwort}

\begin{multicols*}{2}
\textit{Tirakans Reiche} ist ein Rollenspiel.

Diese Regeln und alle zugehörigen Materialien reichen aus, um Abenteuer in der Welt Tirakan zu spielen.
Du kannst diese Regeln frei verwenden und mit Freunden Abenteuer in der Welt bestreiten.

\begin{pquote}{Johann W. von Goethe, Faust}
Ihr wißt, auf unsren deutschen Bühnen\\
Probiert ein jeder, was er mag;\\
Drum schonet mir an diesem Tag\\
Prospekte nicht und nicht Maschinen\\
Gebraucht das groß' und kleine Himmelslicht,\\
Die Sterne dürfet ihr verschwenden;\\
An Wasser, Feuer, Felsenwänden,\\
An Tier und Vögeln fehlt es nicht.\\
So schreitet in dem engen Bretterhaus\\
Den ganzen Kreis der Schöpfung aus\\
Und wandelt mit bedächt'ger Schnelle\\
Vom Himmel durch die Welt zur Hölle.
\end{pquote}

\section{Was ist Rollenspiel?}

Um genau zu sein, reden wir hier von Pen and Paper Rollenspiel, also keinem Computerspiel. Traditionell spielt man ein Rollenspiel mit 2-4 \textit{Spielern} und einer \textit{Spielleiterin} am Tisch, wobei gegenüber einem Brettspiel das Spielbrett durch Charakterbögen und Würfel ersetzt werden. Das Rollenspiel erzählt immer eine Geschichte, die von allen Mitspielern getragen und weitergesponnen wird.

Die Spieler erschaffen für eine Rollenspielrunde Charaktere, also fiktive Figuren, welche auf Charakterbögen festgehalten werden. Der Charakterbogen enthält die Beschreibung des Charakters, seine Herkunft, seine Interessen und seine Fähigkeiten. Letzte sind in Zahlenwerten festgehalten, denn das Handeln im Rollenspiel erfordert sogenannte \textit{Proben} oder \textit{Würfe}, welche den Ausgang einer Aktion bestimmen.

\begin{pexample}[Beispiel]
Tom hat sich entschlossen, in einer Rollenspielgruppe mit Spielleiterin Julia mitzuspielen. Julia hat für ihre Runde ausgesucht, dass es ein Abenteuer im Königreich Asgoran des Jahres 322 geben wird.

Tom beschließt, als Charakter einen adeligen Ritter namens \textit{Jamie Thornton} zu erstellen. Tom nimmt sich einen leeren Charakterbogen und verteilt die verfügbaren Attributs- und Fertigkeitspunkte so, dass Jamie eine 3 in ``Körper'', jeweils eine 2 in ``Wille'' und ``Geschick'' hat, und zudem eine 3 in ``Kampf'' hat. Er beschließt, dass Jamie im Jahr 301 geboren ist, also noch recht jung ist.

Zudem legt Tom für Jamie die Abstammung ``Asgoran'' fest, und wählt als Weg ``Ritter''. Als Bindung wählt er einen Asgoraner Ritterorden.
\end{pexample}

Näheres zu der Erschaffung eines Charakters findest du in \cref{ch:game-concept} und \cref{ch:charaktere}.

Während die Spieler jeweils einen Charakter für das Spiel erstellen, bereitet die Spielleiterin eine Geschichte vor. Diese wird oft auch \textit{Abenteuer} oder \textit{Plot} genannt. Diese Geschichte ist nicht, wie in einem Roman, bis ins Letzte ausformuliert. Stattdessen ist es ein grobes Script, bestehend aus einem möglichen Ablauf, der Beschreibung von Orten, sowie sogenannten Nichtspielercharakteren (auch NSC oder NPC).

Sobald das Spiel beginnt, agieren alle in der Rolle ihrer Charaktere. Die Spielleiterin beschreibt Situationen mit Worten, Karten oder Zeichnungen. Die Spieler sprechen für ihre Charaktere in der ersten Person (``Ich schleiche die Treppe hinauf.''). Kommt es zu Handlungen der Charaktere mit ungewissem Ausgang, so wird auf Proben zurückgegriffen, und es wird gewürfelt.

\begin{pexample}[Beispiel]
Nachdem die Vorbereitungen abgeschlossen sind, trifft sich die Gruppe um Spielleiterin Julia. Dann geht es endlich los, und Julia beschreibt die erste Szene.
\end{pexample}

\begin{pexample}
Julia: ``Es ist der 2. Nebelmond des Jahres 322. Ihr befindet euch in einer Taverne in dem beschaulichen Städtchen Lindfield im Süden Asgorans. Es ist nun bereits spät am Abend, und vor der Tür hat der leichte Nieselregen dafür gesorgt, dass die Schneedecke der letzten Tage von einer dünnen Eisschicht bedeckt ist. Es wird kalt heute Nacht, und glatt. Die Taverne ist gut gefüllt. Mit einem Knarren öffnet sich die Eingangstür, und eine Wolke feinen Regens kommt in den Pub. Direkt gefolgt von einer Gestalt in einer viel zu kleinen Teerjacke.''
\end{pexample}

\begin{pexample}
Dies ist also der Auftakt, und Tom entschließt sich, dass sein Charakter Jamie einmal einen Blick auf den Neuankömmling werfen möchte. Er sagt die Handlungen für Jamie an:

Tom: ``Ich betrachte den Fremden da mal sehr genau, diese schlecht sitzende Jacke ist mir ja schon aufgefallen.''

Julia: ``Du erkennst, dass unter der Kapuze nasse, schwarze Haare in das Gesicht eines alten Mannes fallen. Wirf doch einmal \textit{Wahrnehmung} mit \textit{Instinkt}, um mehr zu erkennen.''
\end{pexample}

Es geht also um ein kooperatives Weiterentwickeln der Geschichte durch das Handeln der Charaktere. Die Spielleiterin hat einen Plan, wohin sich die Geschichte entwickeln könnte, welche Personen auftreten können, und was diese eigentlich für eine Motivation haben. Irgendein Geschehen entwickelt sich um die Charaktere der Spieler herum, und sie werden mit in diese Handlung hinein gezogen.

Wohin diese Geschichte führt, ist ungewiss. Es mag sein, dass etwas Schlimmes verhindert wird, oder dass ein Geheimnis aufgedeckt wird. Die Spielleiterin hat einen groben Plan, die Spieler bestimmen jedoch hauptsächlich den Fortgang.

\subsection{Es geht um das Erzählen}

Denkt man an Computerrollenspiele, so ist die strategische Weiterentwicklung des Charakters der wichtigste Punkt. Er muss zukünftige Kämpfe bestehen können und möglichst optimale Werte für mögliche Herausforderungen haben. Im Pen and Paper Rollenspiel geht es um den Fortgang der Geschichte, um gemeinsame Erlebnisse und Erinnerungen. Das möglichst optimale Ausrichten auf ``starke'' Eigenschaften (sogenanntes \textit{Power-Gaming}) sollte hier nicht im Vordergrund stehen. Dadurch, dass die Geschichte immer gemeinsam weitergetragen wird, gibt es für alle Herausforderungen sehr flexible Lösungen.

Ebenso kommt hier die alte Rollenspielregel zum Tragen, dass das \textbf{Wort der Spielleiterin immer mehr wiegt als die Regeln}. Natürlich sollte es der Normalfall sein, dass die Regeln so angewendet werden, wie sie geschrieben sind, denn es geht für die Spieler auch um einen Rahmen, auf den sie sich verlassen können. Kommt es jedoch zu einer unklaren Regelsituation, so entscheidet die Ansage der Spielleiterin.

\subsection{Kampf im Rollenspiel}

Auch wenn der Fokus im Pen and Paper Rollenspiel weniger auf der Auseinandersetzung mit Fäusten oder Waffen liegt, spielt der Kampf doch eine wichtige Rolle. Nicht jede Situation lässt sich ohne Waffengewalt lösen. Schnell kommt es zu einer Schlägerei, oder die Charaktere planen, einen Händler zu überfallen.

Kampf im Rollenspiel wird anders behandelt als das freie Spiel. Die Zeit läuft in Runden ab, und man veranschaulicht die Situation auf einer Karte auf dem Tisch. Spieler agieren nacheinander, die Spielleiterin steuert ihre NSC. Wunden zeigen an, wie gut es den Charakteren noch geht. Näheres zum Ablauf des Kampfs findest du in \cref{ch:konflikte}.

Im Spiel sollte sich beides die Waage halten. Es mag auch Abenteuer geben, die nur aus einem Kampf bestehen, Tirakans Reiche ist jedoch keine realistische Kampfsimulation. Hier geht es darum, einen Konflikt möglichst unterhaltsam, cineastisch, und trotzdem spannend auszutragen.

Bei einem Kampf im \textit{Tirakans Reiche} System sollten aufgrund der Besonderheiten (Reaktionen, Aktionen stehlen etc.) allerdings immer folgende Dinge eingehalten werden:

\begin{itemize}
    \item Verwendet immer eine Karte. Eine Grundrisskarte der Situation sorgt dafür, dass es keine Missverständnisse bei der Positionierung gibt, egal wie kurz der Kampf ist. Eine Karte kann eine vorgefertigte, aufwändige Karte, aber auch ein schnell gezeichneter Grundriss sein.
    \item Verwendet immer einen Massstab. Die Charaktere haben unterschiedliche Bewegungsreichweiten, und bei einigen ist es auch ihre Stärke, dass sie sich besonders weit bewegen können. Diesen Charakteren und Gegnern wird ohne Massstab ihre Stärke genommen.
    \item Benutzt einen Initiative-Tracker. Initiative bestimmt im Kampf die Reihenfolge der Beteiligten. Initiative-Tracker bedeutet, dass diese Reihenfolge für alle sichtbar aufgeschrieben ist. Es ist bei Tirakans Reiche für die Spieler wichtig zu wissen, wann sie selbst wieder an der Reihe sind (denn dann verfallen u.A. aufgehobene eigene Handlungen).
\end{itemize}

\section{Die Welt Tirakan}
TBD

\end{multicols*}
